% !PRINTED 477
\weberpage{477}
Wendet man darauf die Substitutionen $\omega$ und $\tau$ an, so sieht man, dass $v$ ungeändert bleibt und daher alle Permutationen der Gruppe $P_{168}$ gestattet.

Nun ist $P_{168}$ Theiler der symmetrischen Permutationsgruppe von sieben Elementen vom Index 30 und Theiler der alterniren\-den Gruppe vom Index 15. Folglich ist $v$ Wurzel einer Gleichung 30\,sten Grades, deren Co\"efficienten symmetrische Functionen der $\lambda$ sind, und Wurzel einer Gleichung 15\,ten Grades, deren Co\"efficienten noch das Differenzenproduct der $\lambda$ enthalten. Sind die $\lambda$ die Wurzeln einer Gleichung 7\,ten Grades ohne Affect, so ist $v$ die Wurzel einer Resolvente 30\,sten Grades, die durch Adjunction der Quadratwurzel aus der Discriminante in zwei Factoren 15\,ten Grades zerfällt.

Wenn ausser den symmetrischen Functionen der $\lambda$ die Grösse $v$ dem Rationalitätsbereiche angehört, sei es, dass sie von vornherein rational ist, oder dass der Rationalitätsbereich durch Adjunction von $v$ erweitert wird, so sind die $\lambda$ die Wurzeln einer speciellen Gleichung 7\,ten Grades, deren \gesperrt{Galois}'sche Gruppe vom 168\,sten Grade ist. Was wir nun noch zu beweisen haben, ist, dass sich diese specielle Art von Gleichungen 7\,ten Grades auf das Formenproblem der Gruppe $G_{168}$ zurückführen lässt.

Um dies zu erreichen, müssen wir drei rationale Functionen $X_1$, $X_2$, $X_3$ der Wurzeln $\lambda$ zu bilden suchen, die, wenn die Permutationen der Gruppe $P_{168}$ ausgeführt werden, die entsprechenden linearen Substitutionen der Gruppe $G_{168}$ erfahren, d.\,h. die, wenn $\pi$ eine Permutation aus $P_{168}$ und $A$ die entsprechende Substitution aus $G_{168}$ ist, durch Ausführung der Permutation $\pi$ in $A(X_1, X_2, X_3)$ übergehen.

Setzen wir diese Functionen $X_1, X_2, X_3$ für die Variablen in die Invarianten der Gruppe $G_{168}$ ein, so gehen diese Invarianten in Functionen der $\lambda_i$ über, die durch die Permutationen der Gruppe $P_{168}$ ungeändert bleiben, und die folglich dem Rationalitätsbereiche angehören. Die Berechnung der Werthe der Functionen $X_1, X_2, X_3$ aus diesen Werthen der Invarianten ist dann das Formenproblem für die $G_{168}$. Irgend eine durch die Substitutionen von $G_{168}$ veränderte Function der $X_1, X_2, X_3$ ist Wurzel einer Resolvente der gegebenen Gleichung 7\,ten Grades, und zwar, da die Gruppen $G_{168}$ und $P_{168}$ einfach sind, eine Totalresolvente.
